% =============================================================
%  Recovery sequence — the restart attempts, how their outcomes are
%  classified, the steps of one restart, the restart at reduced
%  voltage and the power compensation by a parallel link. Principles;
%  the specified times and the recorded durations of one project are
%  in pugalur-thrissur.tex.
% =============================================================
\subsection{Recovery Sequence}
\label{sec:dcf-sequence}

\subsubsection{Restart attempts}
\label{sec:dcf-sequence-attempts}

A recovery sequence is a series of up to three restart attempts, each with
its own target state~[24], [25]:

\begin{enumerate}[itemsep=0.2em]
  \item \textbf{First restart} at the pre-fault direct voltage and power
        transfer. Its time limit covers fault clearing, discharge and
        deionisation, and is set separately for pole-to-ground and
        pole-to-pole faults, the latter longer because of the current decay
        of Section~\ref{sec:dcf-faults-types}.
  \item \textbf{Second restart} at a reduced direct voltage, with the power
        curtailed accordingly. Its duration is governed by the on-load
        tap-changer and the recharging of the switchgear drives,
        Section~\ref{sec:dcf-sequence-reduced}, and is specified only as
        slow.
  \item \textbf{Changeover to STATCOM operation} if the second restart fails:
        the DC line is de-energised permanently and both converter stations
        reconfigure to reactive power support of the AC network,
        Section~\ref{sec:dcf-zones-statcom}.
\end{enumerate}

The STATCOM requirement comes from the AC network. The loss of the
transmission changes the load flows and raises the need for reactive power,
and a VSC supplies it independently of the active power, so no additional
FACTS equipment is needed for the contingency. The same requirement applies
after a permanent fault in the DC cable or in one of the converter
stations~[24].

\subsubsection{Classification of faults and attempts}
\label{sec:dcf-sequence-classification}

The outcome of each attempt is classified by the passive charging of the DC
circuit and by a waiting time after a successful restart~[24]. A fault
detected anywhere in the link during a restart attempt makes the fault
\emph{permanent}, and no further attempt follows. An attempt is
\emph{successful} when the power is restored; a fault within a set interval
after that point is \emph{consecutive} and starts the next attempt of the
same sequence. When the interval passes without a fault the sequence is
complete and the next fault opens a new one with all attempts. The
specification also fixes how many complete sequences the equipment has to
withstand within a longer period, which sets the thermal ratings of
Section~\ref{sec:dcf-ratings-thermal}.

A sequence ends unsuccessfully either by a permanent fault or by a third
consecutive fault. In both cases both stations change over to STATCOM
operation after the fault clearing.

\subsubsection{Steps of one restart}
\label{sec:dcf-sequence-steps}

The first restart after a pole-to-ground fault runs in eight steps,
coordinated between the two converter stations and, on a hybrid line, the
transition station over the HVDC telecontrol; without telecontrol the
sequence cannot start~[25].

\begin{enumerate}[itemsep=0.2em]
  \item \textbf{Fault isolation.} The converters block on fault detection
        and the AC circuit breakers of both stations open. Until they open,
        the fault current flows through the diode path.
  \item \textbf{Fault current decay.} With the breakers open the fault
        current decays with the $L/R$ time constant of the loop, set mainly
        by the star-point units. The overvoltage of the healthy pole, limited
        by the arresters, discharges slightly.
  \item \textbf{Discharge and deionisation.} The discharge unit of the lead
        station closes and removes the charge of the healthy pole; the
        bypass switch of the insertion resistor opens. The arc at the fault
        location deionises.
  \item \textbf{Condition verification.} Both stations check breaker and
        bypass-switch positions, arm, DC and star-point currents and DC
        voltages against their thresholds.
  \item \textbf{Tap-changer adjustment.} Only for a restart at reduced
        voltage: the tap-changers move to the start position given by the
        sequence counter.
  \item \textbf{Restart initiation.} The restart release closes the AC
        circuit breaker of the rectifier station first, then of the inverter
        station.
  \item \textbf{Passive charging.} The DC circuit charges through the
        insertion resistor and the diodes. The build-up of the voltages and
        the decay of the charging current confirm the integrity of the DC
        circuit before the resistor is bypassed or the converter is
        deblocked~[24]. Since the sub-module capacitors are still charged,
        this step is shorter than a cold start.
  \item \textbf{Converter deblocking.} The converters deblock and ramp to the
        reference direct voltage. The rectifier station restores its reactive
        power first; the inverter station then restores active and reactive
        power to the pre-fault values.
\end{enumerate}

The restart after a pole-to-pole fault omits the discharge step, since there
is no trapped charge, but has to wait for the freewheeling current to reach
zero and for the deionisation time after it before the circuit breakers
reclose~[24]. Passive charging is also the test for a permanent fault: on a
persisting pole-to-ground fault the direct voltage builds up in one pole
only after the circuit breaker closes, the current stays limited by the
insertion resistor, the protection detects the condition, the start-up is
aborted and the healthy pole is discharged~[24].

\subsubsection{Restart at reduced direct voltage}
\label{sec:dcf-sequence-reduced}

A half-bridge MMC has a narrow direct-voltage range: to avoid
over-modulation, the direct voltage to ground has to stay above the peak of
the AC voltage the converter produces at its terminals, so the AC-side
voltage has to be lowered with the direct voltage. The ratio of the
interface transformer is changed with its on-load tap-changer, before the
circuit breaker recloses, since the link has to recover to the reduced
voltage immediately after deblocking. Depending on the number of taps this
takes up to a minute~[24]. The operator preselects the reduced voltage
within a range the design allows, and a link already running below the
preselected level is kept from being raised to it by the recovery~[24]. In
every case the sequence has to respect the charging state of the drives of
the breakers and fast switches, that is of the discharge switch and the
insertion-resistor bypass switch~[24].

\subsubsection{Power compensation by a parallel link}
\label{sec:dcf-sequence-compensation}

Where a project consists of two independent links, the specification can
require the active power of the faulted link to be transferred to the
healthy one, as far as its rating allows, within a few tens of
milliseconds~[24]. The healthy link then carries the power until the faulted
link has deblocked, and the original sharing is restored afterwards with a
long time constant, so that the AC network sees no second disturbance; this
slow return applies only to the share the healthy link took over, and a
share it could not take over is restored immediately after
deblocking~[24]. A scheme with a single link has no such path: the loss of
transmission until deblocking and the ramp to pre-fault power act on the AC
network at both ends, and other measures to limit that impact become a
requirement of the project~[24].

\FloatBarrier
